\section{Terminal pullback and exact reciprocal sum}
\label{sec:terminalization}

Only pullback stability of surjections in $\mathbf{Set}$ remains.  For $X\in\mathcal H_B$ put
\[
T_X=\operatorname{Tor}_{E_X}(\gamma_X)\times I_X,
\qquad
\mathcal T=\coprod_{X\in\mathcal H_B}T_X.
\]
A realization $r\in\mathcal R_X$ contains a labelled correspondence $R_X$ covering $T_X$.  Let
\[
\mathcal E
=
\{(X,r,t,e):X\in\mathcal H_B,\ r\in\mathcal R_X,\ t\in T_X,
\ e\text{ is an edge of }R_X\text{ ending at }t\}.
\]
There is a natural surjection
\begin{equation}\label{eq:universal-edge-surjection}
\mathcal E
\twoheadrightarrow
\mathcal R\times_{\mathcal H_B}\mathcal T.
\end{equation}
Every interval configuration labelling an edge of $\mathcal E$ has reciprocal value below $2$ by \eqref{eq:physical-W-less-2}.

Fix a sufficiently large integer $k$.  Theorem~\ref{thm:scalar-terminality} gives the nonempty set $\mathcal X_k$ from \eqref{eq:term-k}.  Form the pullback
\[
\begin{CD}
\mathcal R_k @>>> \mathcal R\\
@VVV @VVV\\
\mathcal X_k @>>> \mathcal H_B.
\end{CD}
\tag{8.1}\label{eq:real-k-square}
\]
Since $\mathcal R\to\mathcal H_B$ is surjective, so is
\begin{equation}\label{eq:real-k}
\mathcal R_k\twoheadrightarrow\mathcal X_k.
\end{equation}

For $X\in\mathcal X_k$ we have $\gamma_X=0$ and $k\in I_X$, hence there is a canonical point
\[
(0,k)\in T_X.
\]
These points form a section
\[
s_k:\mathcal X_k\longrightarrow\mathcal T
\]
over $\mathcal H_B$.  Combining it with \eqref{eq:real-k-square} gives a map
\[
\mathcal R_k\longrightarrow
\mathcal R\times_{\mathcal H_B}\mathcal T.
\]
Pulling \eqref{eq:universal-edge-surjection} back along this map gives the second cartesian square
\[
\begin{CD}
\mathcal E_k @>>> \mathcal E\\
@VVV @VVV\\
\mathcal R_k @>>> \mathcal R\times_{\mathcal H_B}\mathcal T.
\end{CD}
\tag{8.2}\label{eq:terminal-edge-pullback}
\]
Thus $\mathcal E_k\to\mathcal R_k$ is surjective.  Since $\mathcal X_k$ is nonempty, both $\mathcal R_k$ and $\mathcal E_k$ are nonempty.  Any element of $\mathcal E_k$ supplies an interval configuration $S\in\mathcal C$ satisfying
\begin{equation}\label{eq:terminal-state}
\overline W(S)=0,
\qquad
\nu(S)=k,
\qquad
W(S)<2.
\end{equation}
No coherent choice as $k$ varies is required.

Let
\[
q_{\mathbf Z}:\mathbf Q\longrightarrow\mathbf Q/\mathbf Z
\]
be the quotient map.  By definition of $\overline W$, the square
\[
\begin{CD}
\mathcal C @>{W}>> \mathbf Q\\
@V{\overline W}VV @VV{q_{\mathbf Z}}V\\
\mathbf Q/\mathbf Z @= \mathbf Q/\mathbf Z
\end{CD}
\tag{8.3}\label{eq:value-residue-square}
\]
commutes.  Hence $\overline W(S)=0$ implies
\[
W(S)\in\ker q_{\mathbf Z}=\mathbf Z.
\]
Since $\nu(S)=k>0$, the configuration is nonempty and $W(S)>0$.  Together with \eqref{eq:terminal-state},
\[
0<W(S)<2,
\qquad
W(S)\in\mathbf Z,
\]
so $W(S)=1$.  The $k$ connected components of $S$ are pairwise nonadjacent intervals of cardinality $2$ or $3$.  This proves Theorem~\ref{thm:main}.
